\section{Future Directions}
\label{sec:future}

The functional taxonomy exposes several gaps that are difficult to see from architecture-level comparisons alone. The central problem is no longer whether a system can issue several searches, but whether it can maintain a useful information state over long horizons, preserve the dependencies required by later reasoning, and improve without accumulating brittle or stale behavior.

\subsection{Retrieval for Future Utility}
Most retrievers are trained to estimate immediate query--document relevance. Multi-hop systems additionally need evidence whose value is conditional: a bridge passage may not answer the current question, yet may identify the entity, relation, constraint, or vocabulary required for a later search. Future work should therefore distinguish \emph{answer evidence}, \emph{bridge evidence}, and \emph{search-enabling evidence}, and learn objectives that value their different roles.

A key research question is how to estimate future utility without requiring a fully annotated gold chain. Possible signals include successful downstream queries, complete-chain recovery, counterfactual removal of bridge passages, and process rewards tied to subsequent search progress. Evaluation should report whether the method improves entry evidence, later-hop reachability, or both.

\subsection{State Representation and Error Recovery in Long-Horizon Search}
A long trajectory cannot be controlled reliably from an unstructured transcript alone. The system must represent what has been established, what remains uncertain, which sources support each claim, which queries have already failed, and which branches remain open. State representations may combine textual notes, entity--relation structures, query histories, uncertainty estimates, and source provenance.

Long-horizon systems also require recovery mechanisms. One incorrect entity resolution can redirect all later searches, while repeated reformulation can amplify a false assumption. Future methods should detect inconsistent intermediate states, revisit earlier decisions, compare alternative branches, and stop when additional search is unlikely to change the result. These capabilities should be evaluated through controlled perturbations of intermediate evidence rather than only through naturally occurring failures.

\subsection{Joint Optimization without Losing Component Accountability}
Modern systems increasingly optimize retrieval, reasoning, tool use, and generation jointly. End-to-end optimization can improve final task reward, but it can also obscure which component changed and whether the system learned a shortcut. A useful research direction is modular or constrained joint learning in which search, evidence integration, and generation share feedback while retaining stage-level observability.

This requires reporting interfaces between components: the candidate pool before and after search control, the evidence set before and after integration, and the trajectory before and after feedback updates. Without these controls, a gain attributed to retrieval may instead result from a stronger reader, a larger context, or additional model calls.

\subsection{Evidence Preservation with Explicit Provenance}
Compression and note taking should preserve not only semantic content but also the dependency structure needed by later reasoning. Important details include entity identity, negation, numerical values, temporal qualifiers, uncertainty, and the source from which each statement was derived. A compact note that removes provenance or merges conflicting claims may be efficient but unusable for grounded synthesis.

Future systems should represent compressed evidence as structured claims linked to source spans. Compression evaluation should combine token reduction with relation preservation, answerability under controlled evidence, and citation recovery. This would make it possible to separate harmless wording compression from deletion of a reasoning-critical bridge.

\subsection{Faithful Long-Form Synthesis}
Deep research shifts the output from a short answer to a report containing many claims with different evidential status. The system must decide which claims are directly stated by sources, which are derived by combining sources, and which remain speculative. It must also reconcile contradictions, avoid citing a source for a stronger claim than it supports, and preserve temporal context.

A promising direction is claim--source graph construction before report generation. Such a representation could support coverage checks, contradiction detection, citation placement, and targeted re-search for unsupported sections. Evaluation should measure claim-level support and completeness in addition to document-level citation counts or holistic report ratings.

\subsection{Reusable Experience and Skills}
Feedback learning introduces a second horizon: the system must improve across tasks rather than only within one trajectory. Current experience stores and skill libraries often lack clear policies for retention, retrieval, revision, and deletion. A successful trajectory may encode a domain-specific shortcut that transfers poorly, while an old skill may remain highly ranked after the environment changes.

Future work should evaluate experience reuse on genuinely new tasks, report the retrieval accuracy of memories or skills, and measure negative transfer and forgetting. Skill representations should expose their preconditions, expected effects, evidence requirements, and failure history. Versioning and conflict resolution are necessary if independently learned procedures are to be composed safely.

\subsection{Benchmarks for Open and Repeated Information Seeking}
Most multi-hop benchmarks evaluate a fixed corpus and a single episode. Open-web research introduces changing documents, heterogeneous source quality, access cost, and incomplete gold evidence. Feedback learning additionally requires repeated tasks and a temporal split between experience acquisition and later evaluation.

New benchmarks should release or reconstruct action traces, source visits, retrieved evidence, stopping decisions, and costs. They should distinguish static knowledge access from live-web search, and current-task adaptation from improvements caused by prior experience. For long-form outputs, benchmark annotations should connect claims to source spans rather than provide only a global score.

\subsection{Safety and Governance of Autonomous Information Seeking}
Adaptive retrieval expands the attack surface of a language model. Malicious pages can manipulate queries or tool calls; biased source selection can systematically omit perspectives; persistent memory can retain private or poisoned information; and long-running agents can incur unbounded cost. These risks arise during acquisition and learning, not only in the final text.

Safety evaluation should therefore include source trust, prompt-injection resistance, privacy controls, memory deletion, budget enforcement, and auditability of search actions. Persistent feedback mechanisms require particular scrutiny because a single poisoned trajectory can affect many future tasks. The same process-level logging proposed for scientific evaluation is also necessary for governance and incident analysis.
